\documentclass[pdftex,12pt,a4paper]{report}
\usepackage[utf8]{inputenc}
\usepackage[english]{babel}
\usepackage{natbib}
\usepackage{graphicx}
\usepackage{subfiles}
\usepackage{amsmath}
\usepackage{amsthm}
\usepackage{tikz}
\usepackage{fancyhdr}
\usepackage{float}
\usepackage{algorithm}% http://ctan.org/pkg/algorithms
\usepackage{algpseudocode}% http://ctan.org/pkg/algorithmicx
\usepackage{wrapfig}
\usepackage{lscape}
\usepackage{rotating}
\usepackage{epstopdf}
\usepackage[font=small,labelfont=bf]{caption}
\usetikzlibrary{shapes.geometric, arrows}
\theoremstyle{definition}
\newtheorem{definition}{Definition}[section]
\newtheorem{theorem}{Theorem}[section]
\newtheorem{lemma}[theorem]{Lemma}
 
\theoremstyle{remark}
\newtheorem*{remark}{Remark}
\usepackage{amssymb}
\usepackage{amsfonts}
\usepackage{mathtools}
\usepackage{graphicx}
\usepackage{titlesec}
\usetikzlibrary{positioning}
\graphicspath{{images/}{../images/}}
\usepackage{geometry}
 \geometry{
 a4paper,
 total={210mm,297mm},
 left=30mm,
 right=30mm,
 top=30mm,
 bottom=35mm,
 }
 \pagestyle{fancy}
\fancyhf{}
\chead{\leftmark}
\cfoot{\thepage}
 \newcommand*{\justifyheading}{\raggedleft}
\titleformat{\chapter}[display]
  {\normalfont\huge\bfseries\justifyheading}{\chaptertitlename\ \thechapter}
  {20pt}{\Huge}
\newcommand{\defeq}{\vcentcolon=}
\newcommand{\eqdef}{=\vcentcolon}
\newcommand*{\V}[1]{\mathbf{#1}}%
\newcommand{\norm}[1]{\left\lVert#1\right\rVert}
\newcommand{\justif}[2]{&{#1}&\text{#2}}
\newcommand{\qedwhite}{\hfill \ensuremath{\Box}}
\newcommand\given[1][]{\:#1\vert\:}
\newcommand{\me}{\mathrm{e}}
\DeclarePairedDelimiterX{\infdivx}[2]{(}{)}{%
  #1\;\delimsize\|\;#2%
}
\newcommand{\Conv}{\mathop{\scalebox{1.5}{\raisebox{-0.2ex}{$\ast$}}}}%
\newcommand{\infdiv}{\infdivx}
\renewcommand{\qed}{\hfill\blacksquare}
%\renewcommand{\baselinestretch}{1.3}


\DeclarePairedDelimiter\abs{\lvert}{\rvert}%

\tikzstyle{startstop} = [rectangle, rounded corners, minimum width=3cm, minimum height=1cm,text centered, draw=black]
\tikzstyle{io} = [trapezium, trapezium left angle=70, trapezium right angle=110, minimum width=3cm, minimum height=1cm, text centered, draw=black]
\tikzstyle{decision} = [diamond, minimum width=3cm, minimum height=1cm, text centered, draw=black]
\tikzstyle{process} = [rectangle, minimum width=3cm, minimum height=1cm, text centered, draw=black]
\tikzstyle{arrow} = [thick,->,>=stealth]
\DeclareMathOperator*{\argmax}{arg\,max}
\DeclareMathOperator*{\argmin}{arg\,min}
 
\title{Learning relational structures from birdsong}
\author{Bernardo Pérez Orozco\\Balliol College\\University of Oxford}
\date{ September 1st 2015 }
\begin{document}
 \begin{titlepage} 
\begin{center}

% Upper part of the page. The '~' is needed because \\ % only works if a paragraph has started. \includegraphics[width=0.15\textwidth]{./logo}~\\[1cm]

\textsc{\LARGE Learning relational structures from birdsong}\\[2cm]
\includegraphics[width=0.35\textwidth]{./images/logo.png}\\[2cm]
% Author and supervisor 

{
\large 
\textbf{Bernardo Pérez Orozco}\\
\textsl{Balliol College\\
University of Oxford\\}
\textbf{Supervised by:}\\
\textsl{Prof. Stephen Roberts \& Prof. Thomas Melham}
}
\vfill

% Bottom of the page 
{\large 
\textsl{Dissertation submitted for the degree of}\\
\textsc{Master of Science in Computer Science}\\
\textsl{September 1st 2015}
}

\end{center} 
\end{titlepage}

\begin{abstract}
In this work, we build a hierarchical relational structure for bird species birdsong. In order to do so, we characterise each bird species as a statistical model (using a Variational Bayes Hidden Markov Model) and run Agglomerative Hierarchical Clustering over the pairwise similarity matrix. Each statistical model summarises the first-formant trajectories (feature vectors) of each bird species. Formant trajectories allow us to model the changing resonances of the vocal tract through time.
\par We use a symmetrised KL-Divergence, the Hellinger distance and a model based distance given by comparing the Dirichlet distributions associated with the transition models. We infer the phylo-acoustic trees generated by using different statistical models and metrics and show that there is a clear hierarchical community structure in all cases. Additionally, we also give empirical evidence that the first-formant trajectory is relatively more structured in average for any other formant trajectory of the same birdsong recording, hence setting a precedent of formant trajectories applied as a birdsong feature extraction technique.
\end{abstract}

\renewcommand{\abstractname}{Acknowledgements}
\begin{abstract}
First and foremost, I would like to express my most sincere gratitude to Prof. Steve Roberts, whose patience and guidance proved invaluable during these months. In the same way, I would like to thank to Prof. Tom Melham, whose constant support and advice led me to get the most out of the Balliol experience this year.
\par Thanks are also due to Dr. Davide Zilli for all the discussions, the support and for always lending me a hand when I needed it. I would also like to thank Dr. Carlos Ramírez and Mayra González, for teaching me to reach higher goals.
\par I would also like to thank the National Council for Science and Technology in Mexico (CONACYT) and the Council for Science and Technology in the State of Querétaro (CONCYTEQ) for their financial support throughout the year.
\par Likewise, I cannot stress enough how thankful I am to my parents and my sister, for their endless support in each and all dimensions of my development. Thank you for being so patient, comprehensive and caring. 
\par My deepest gratitude also goes to Rodrigo, for being the kindest and best of friends. To Rebeca, for her endless patience and selflessness. To Cristina, for being with me from the absolute beginning. To Aziza, Siggi, Maximilian, Jennifer and Xetanasha, for opening my eyes to a whole new world beyond my thoughts. And to each and all of my great friends who have pushed me to give my very best: Aarón, Alan, Ari, Andy, Hugo, Meli, Alex, Julia, Max, Polo, Noé, Carlos, Angélica, Angela, César, Rodyka and Marcela. I keep you all in my heart.

\end{abstract}

\tableofcontents
\newpage
\chapter{Introduction}
\subfile{chapters/ch1.tex}
 
\chapter{State of the art}

\subfile{chapters/ch2.tex}

\chapter{Feature extraction and signal analysis}
 
\subfile{chapters/ch3.tex}
 
\chapter{Similarity measures and relational structures}

\subfile{chapters/ch4.tex}

\chapter{Implementation and results}

\subfile{chapters/ch5.tex}

\chapter{Conclusions and future work}

\subfile{chapters/ch6.tex}

\addcontentsline{toc}{chapter}{Bibliography}
\bibliographystyle{unsrt}
\bibliography{bib}


\appendix

\chapter{Metrics review}

\subfile{chapters/app2.tex}

\chapter{Dirichlet and Wishart distributions}

\subfile{chapters/app.tex}

\chapter{Closed forms for KL Divergence and Hellinger distance}

\subfile{chapters/app1.tex}



\end{document}